\chapter{Concurrency revisited}
\label{CH:LOCK2}

Simultaneously obtaining good parallel performance, correctness despite
concurrency, and understandable code is a big challenge in kernel design.
Straightforward use of locks is the best path to correctness, but is not
always possible. This chapter highlights examples in which xv6++
is forced to use locks in an involved way, and examples where xv6++
uses lock-like techniques but not locks.

\section{Locking patterns}

Cached items are often a challenge to lock.
For example, the file system's block cache stores copies of up to
\lstinline{NBUF} disk blocks (see \texttt{kernel/param.h}).
It's vital that a given disk block have at most one copy in the cache;
otherwise, different threads might make conflicting changes to different
copies of what ought to be the same block.

In xv6++ each cached block is stored in a \lstinline{struct buf}
(see \texttt{kernel/buffer\_cache.h}).
A \lstinline{buf} has an embedded sleep-lock (\lstinline{buf::lock})
which helps ensure that only one kernel thread uses a given disk block at a time.
However, that per-buffer lock is not enough: what if a block is not present in
the cache at all, and two threads want to use it at the same time?
There is no existing locked \lstinline{buf} yet (since the block isn't cached),
and thus there is nothing per-block to lock.

Xv6++ deals with this situation by associating an additional lock with the
\emph{set of identities of cached blocks}. In xv6++ this is the
spinlock inside \lstinline{buffer\_cache} (see \texttt{kernel/buffer\_cache.h}).
Code that needs to check whether a block is cached, or that needs to change the
set of cached blocks, must hold the cache's lock. After that code has found (or
installed) the \lstinline{buf} it needs, it can release the cache lock and
then rely on just the per-buffer sleep-lock to protect the block's contents.
This is a common pattern: one lock for the set of items, plus one lock per item.

Ordinarily the same function that acquires a lock will release it.
But a more precise way to view things is that a lock is acquired at the start of
a sequence that must appear atomic, and released when that sequence ends.
If the sequence starts and ends in different functions, or different threads,
or on different CPUs, then the lock acquire and release must do the same.
The function of the lock is to force other uses to wait, not to pin a piece of
data to a particular agent.

In xv6++ there are many cases where a kernel thread sleeps while holding a
sleep-lock and later continues on a different CPU. For example, file system
code often holds an inode sleep-lock across disk I/O. The code can block waiting
for the disk, and may resume execution on a different CPU. Thus the lock may be
acquired and released on different CPUs; this is safe because the lock's purpose
is mutual exclusion, not CPU affinity.

Freeing an object protected by a lock embedded in the object is delicate, since
owning the lock is not enough to guarantee that freeing is correct.
The problem case arises when some other thread is waiting in \lstinline{acquire}
to use the object; freeing the object implicitly frees the embedded lock, which
will cause the waiting thread to malfunction.
One solution is to track how many references to the object exist, so that it is
only freed when the last reference disappears.
Xv6++ uses this strategy in many places: for example, \lstinline{file} objects
have a reference count (\lstinline{file::ref}), inodes have a reference count
(\lstinline{inode::ref}), and buffers have \lstinline{buf::refcnt}.
Pipes also track whether there are still open endpoints (see
\texttt{kernel/pipe.h} and \texttt{kernel/pipe.cpp}), so the pipe is only torn
down when both ends are closed.

Usually one sees locks around sequences of reads and writes to sets of related
items; the locks ensure that other threads see only completed sequences of
updates (as long as they, too, lock). What about situations where the update is
a simple write to a single shared variable?

In xv6++ there are places where a small shared flag is protected by a lock to
avoid a data race. For example, process state such as \lstinline{process::killed}
is protected by \lstinline{process::lock} when read or written.
If there were no lock, one thread could write a flag at the same time that
another thread reads it. This is a \indextextx{race}, and in C/C++ a race yields
\indextext{undefined behavior}\footnote{See ``Data races'' in
\url{https://en.cppreference.com/w/cpp/language/memory_model}}.
The locks prevent the race and avoid undefined behavior.

One reason races can break programs is that, without locks or equivalent
constructs, the compiler may generate machine code that reads and writes memory
in ways quite different than the source code. For example, the machine code for
a thread reading a boolean flag could cache it in a register and never re-load
it, meaning it might never observe another thread's writes. Locks (and the
atomic operations they entail) prevent such incorrect caching.

\section{Lock-like patterns}

In many places xv6++ uses a reference count or a flag in a lock-like way to
indicate that an object is allocated and should not be freed or re-used.
A process's state acts this way, as do the reference counts in
\lstinline{file}, \lstinline{inode}, and \lstinline{buf} objects.
While in each case a lock protects the flag or reference count itself,
it is the latter that prevents the object from being prematurely freed.

The file system uses inode reference counts as a kind of shared lock that can
be held by multiple threads, in order to avoid deadlocks that would occur if
the code used only ordinary locks.
For example, pathname lookup evaluates one directory component at a time.
The lookup code must not hold multiple inode sleep-locks at once, since it could
deadlock with itself if the pathname includes a dot (e.g., \lstinline{a/./b}).
It might also deadlock with a concurrent lookup involving the directory and
\lstinline{..}. As Chapter~\ref{CH:FS} explains, the solution is to carry the
next inode to the following iteration with its reference count incremented
(using \lstinline{idup()} / \lstinline{iput()} in xv6++'s \lstinline{file_system}
interface), but not locked.

Some data items are protected by different mechanisms at different times, and
may at times be protected from concurrent access implicitly by the structure of
the code rather than by explicit locks.
For example, when a physical page is free it is protected by the allocator's
lock (see \texttt{kernel/page\_allocator.h} and \texttt{kernel/page\_allocator.cpp}).
If the page is then allocated as part of a pipe's storage, it is protected by
the pipe's embedded lock. If the page is re-allocated for a new process's user
memory, it is not protected by a lock at all; instead, the fact that the
allocator will not give that page to any other process (until it is freed)
protects it from concurrent access.

The ownership of a process's memory is complex:
first the parent allocates and manipulates it in \lstinline{fork()},
then the child uses it, and (after the child exits) the parent again
owns related resources and eventually frees them.
There are two lessons here: a data object may be protected from concurrency in
different ways at different points in its lifetime, and the protection may take
the form of implicit structure rather than explicit locks.

A final lock-like example is the need to disable interrupts around code that
depends on per-CPU identity.
In xv6++ the CPU ID is derived from the RISC-V thread pointer register (\lstinline{tp}),
and code that queries ``current CPU'' or ``current process'' must ensure it will
not be preempted and moved in the middle of the sequence.
Xv6++ provides \lstinline{cpu_manager::push_off()} and \lstinline{pop_off()}
to disable interrupts around such sequences (see \texttt{kernel/cpu\_manager.h}).

\section{No locks at all}

There are a few places where xv6++ shares mutable data with no locks at all.
One is in the implementation of spinlocks themselves, although one could view
the RISC-V atomic instructions as relying on locks implemented in hardware.

Xv6++ contains cases in which one CPU or thread writes some data and another CPU
or thread reads it, but there is no specific lock dedicated to protecting that
data. For example, during \lstinline{fork()}, the parent writes the child's user
memory pages, and the child later reads those pages; no explicit lock protects
those pages. This is not strictly a locking problem, since the child does not
start executing user code until after the parent has finished initializing the
child's memory and state.

It is, however, a potential memory-ordering problem (see Chapter~\ref{CH:LOCK}),
since without a memory barrier there's no reason to expect one CPU to see
another CPU's writes. In practice the parent releases locks as it finishes the
setup work and the child acquires locks as it starts running; the memory barriers
in spinlock acquire/release ensure that the child's CPU sees the parent's writes.

\section{Parallelism}

Locking is primarily about suppressing parallelism in the interests of
correctness. Because performance is also important, kernel designers often have
to think about how to use locks in a way that both achieves correctness and
allows parallelism. While xv6++ is not systematically designed for high
performance, it's still worth considering which operations can execute in
parallel, and which might conflict on locks.

Pipes are an example of fairly good parallelism. Each pipe has its own lock, so
different processes can read and write different pipes in parallel on different
CPUs. For a given pipe, however, the writer and reader must wait for each other
to release the lock; they can't read and write the same pipe at the same time.
A read from an empty pipe (or write to a full pipe) must block, but this is not
primarily a consequence of locking.

Context switching is a more complex example. Two kernel threads, each executing
on its own CPU, can call \lstinline{yield()}, \lstinline{sched()}, and
\lstinline{swtch()} at the same time, and their calls will execute in parallel.
Each thread typically holds a per-process lock while transitioning state; since
they are different process locks, they don't have to wait for each other.
Once inside the scheduler loop, however, multiple CPUs may conflict while
searching the process table for runnable processes and while acquiring per-process
locks to examine and change process state.

Another example is concurrent calls to \lstinline{fork()} from different
processes on different CPUs. The calls may have to wait for each other on global
locks (for example the allocator lock when allocating pages, or a PID lock when
assigning PIDs) and on per-process locks needed to locate an unused process
slot. On the other hand, once resources are allocated, the two \lstinline{fork()}
calls can copy user memory pages and build page tables in parallel.

The locking scheme in each of the above examples sacrifices parallel performance
in certain cases. In each case it's possible to obtain more parallelism using a
more elaborate design. Whether it's worthwhile depends on details: how often the
operations run, how long contended locks are held, how many CPUs might run
conflicting operations at the same time, and whether other code paths are the
real bottleneck. It can be difficult to guess whether a given locking scheme
will cause performance problems, so measurement on realistic workloads is often
required.

\section{Exercises}

\begin{enumerate}

\item Modify xv6++'s pipe implementation to allow a read and a write to the same
pipe to proceed in parallel on different CPUs.

\item Modify xv6++'s process scheduler to reduce lock contention when different
CPUs are looking for runnable processes at the same time.

\item Eliminate some of the serialization in xv6++'s \lstinline{fork()}.

\end{enumerate}
